\documentclass{article}

\textwidth=6.5in
\textheight=8.5in
\voffset=-54pt
\oddsidemargin=0in
\evensidemargin=0in
\setlength{\parskip}{8pt}
\setlength{\parindent}{0pt}

\usepackage{workbook}

\usepackage[workshop]{handout}
\renewcommand{\workshopnote}{CA Notes.} 
\usepackage{lipsum}
 \usepackage{microtype}
 \usepackage{vwcol}

\begin{document}
\htitle{Workshop 10: Integration Techniques and Slicing to Find Area}
\begin{workshop}
    Welcome to Workshop 10! The objectives for this workshop are for students to:
    \begin{itemize}
        \item discern between integrands that are suitable for integration by substitution or by parts, and evaluate these integrals.
        \item carry out advanced problem types with integration by parts (including IBP twice and combining substitution and IBP).
        \item finding the area enclosed between curves and/or boundaries via vertical or horizontal slicing to generate a solvable integral.
    \end{itemize}
    Here is a sample lesson plan:
    \begin{itemize}
        \item First 25 minutes: \#1 and \#2 -- These are more straightforward examples of integration by substitution and by parts that students have already encountered frequently. Students can work on these at their tables or at their boards. One purpose of these exercises is to get practice spotting when an integrand is suitable for substitution or parts, so reinforce questions like ``Do you see a function and its derivative within the integrand?" or ``Does it look like there a product in the integrand that can't be simplified with a substitution?".
        \item 15 minutes: \#3 Students might benefit from a few minutes to work at tables, then discuss, via a conversation as a class, the appropriate notation and overall slice-approximate-sum-limit procedure to compute area in these types of problems.
        \item 35 minutes: \#4 Students should practice evaluating the same integral by (1) horizontally slicing and (2) and integrating by parts in groups at the boards. Please review horizontal slicing as a class, as this topic tends to be more challenging for students.
        \item If time permits: Students can begin to work on the extra practice integrals by parts.
    \end{itemize}
\end{workshop}
\begin{enumerate}
\item In each part below you are given a pair of integrals. Use integration by parts to evaluate one of the integrals per pair, and evaluate the other using integration by substitution.
\begin{enumerate}
    \item
        \begin{enumerate}
        \begin{multicols}{2}
        \item $\displaystyle \int xe^{-x}\,dx$
        
        \begin{solution}
        
            Let's use integration by parts.
            
\begin{tabular}{ l l }
  $u=x$ &  $dv = e^{-x}dx$  \\
  $du=dx$ &  $v=-e^{-x}$  \\
\end{tabular}
\begin{align*}
\displaystyle \int xe^{-x}\,dx &= -xe^{-x}-\displaystyle \int -e^{-x}\,dx\\&=xe^x-e^{-x}+C
\end{align*}
        \end{solution}
\columnbreak
        \item $\displaystyle \int xe^{-x^2}\,dx$
        
        \begin{solution}
        Let's use the substitution $u=-x^2$ and $du=-2xdx$, where $-\frac{du}{2}=xdx$.

$\displaystyle -\dfrac{1}{2} \int e^u\,du = -\dfrac{1}{2}e^u+C=-\dfrac{1}{2}e^{-x^2}+C$
        
        \end{solution}
        \end{multicols}
        
        \end{enumerate}
    \pspace{\vfill}
    \item
        \begin{enumerate}
        \begin{multicols}{2}
        \item $\displaystyle \int \dfrac{\ln x}{x^2}\,dx$
        
        \begin{solution}
         Let's use integration by parts. 
            
\begin{tabular}{ l l }
  $u=\ln x$ &  $dv = \frac{1}{x^2}dx$  \\
  $du=\frac{1}{x}dx$ &  $v=-
  \frac{1}{x}$  \\
\end{tabular}
\begin{align*}
\displaystyle \int \dfrac{\ln x}{x^2}\,dx&=-\frac{\ln x}{x}-\displaystyle \int \frac{-1}{x^2}\,dx\\
&=-\frac{\ln x}{x}+\displaystyle \int \frac{1}{x^2}\,dx \\
&=-\frac{\ln x}{x}-\frac{1}{x}+C \\
\end{align*}


        \end{solution}
        \columnbreak
        \item $\displaystyle \int \dfrac{\ln x}{x}\,dx$ 

        \begin{solution}
        Let's use the substitution $u=\ln x$ and $du=\frac{1}{x}dx$.

        $\displaystyle \int u \,du = \frac{1}{2}u^2+C=\frac{1}{2}(\ln x)^2+C$
        \end{solution}
        \end{multicols}
       
        \end{enumerate}
        \pspace{\vfill}
     \pspace{\newpage}   
     \item
        \begin{enumerate}
        \begin{multicols}{2}
        \item $\displaystyle \int x\cos(x^2)\,dx$
        
        \begin{solution}
            Let's use the substitution $u=x^2$ and $du=2xdx$, where $\frac{du}{2}=xdx$.

            $\displaystyle \frac{1}{2} \int \cos u\,du=\frac{1}{2} \sin u +C =\frac{1}{2} \sin (x^2) +C $
        \end{solution}
\columnbreak
        \item $\displaystyle \int x\cos x\,dx$
        
        \begin{solution}
        
        \begin{tabular}{ l l }
  $u=x$ &  $dv = \cos xdx$  \\
  $du=dx$ &  $v=\sin x$  \\
\end{tabular}
\begin{align*}
\displaystyle \int x\cos x\,dx&=x\sin x-\displaystyle \int \sin x\,dx\\
&=x\sin x + \cos x +C\\
\end{align*}
        \end{solution}
        \end{multicols}
        
        \end{enumerate}
        \end{enumerate}

    \pspace{\vfill}
    
    \item Evaluate the following integrals. (Hint: You should be able to make use of your prior work in \#1c(ii).)
    \begin{enumerate}
    \item $\displaystyle \int x^3\cos (x^2)\, dx$ 
    
    \begin{solution}
        Let $u=x^2$ and $du=2xdx$, where $\frac{du}{2}=xdx$.
        
        We can rewrite the integral as $\displaystyle \int x^2\cdot x\cdot\cos (x^2)\, dx$. According to our substitution, 
        \begin{center}
        $\displaystyle \int x^2\cdot x\cdot\cos (x^2)\, dx$ = $\displaystyle \frac{1}{2}\int u \cos (u)\, du$
        \end{center}

        According to our work in \#1c(ii), $\displaystyle \frac{1}{2} \int u\cos u\,du = \frac{1}{2}(u\sin u + \cos u) +C$. 
        
        Since $u=x^2$,
        \begin{center}
            $\displaystyle \int x^3\cos (x^2)\, dx = \frac{1}{2}[x^2\sin (x^2) + \cos (x^2)]+C$
        \end{center}
        \end{solution}
        \pspace{\vfill}
  
    \item $\displaystyle \int x^2\sin (x)\, dx$ 
   
   
        \begin{solution}
        Let's use integration by parts. 

                  
\begin{tabular}{ l l }
  $u=x^2$ &  $dv = \sin xdx$  \\
  $du=2xdx$ &  $v=-\cos x$  \\
\end{tabular}
\begin{align*}
\displaystyle \int x^2\sin(x)\,dx&=-x^2\cos x-\displaystyle \int -2x\cos x\,dx\\
&=-x^2\cos x +\displaystyle 2 \int x\cos x\,dx\\
\end{align*}

We can use our work from \#1c(ii) to finish evaluating this integral. 

\begin{align*}
\displaystyle \int x^2\sin(x)\,dx&=-x^2\cos x +\displaystyle 2 \int x\cos x\,dx \text{\hspace{0.5in}(copied from above)}\\
&=-x^2\cos x +2(x\sin x+\cos x)+C\\
&=-x^2\cos x+2x\sin x +2\cos x +C
\end{align*}

        \end{solution}
        
 \pspace{\vfill}
 \end{enumerate}
        
        
  
     
       
       
       

\pspace{\newpage}
\item Let $\mathcal{R}$ be the region enclosed by $y=3\sqrt{x}$ and $y=x$.
\begin{enumerate}
    \item Sketch the region $\mathcal{R}$ below.
    \probonly{
   \begin{center}
    \begin{tikzpicture}
        \begin{axis}[xmin=0, xmax = 10, xtick={0,1,2,3,4,5,6,7,8,9}, ytick={0,1,2,3,4,5,6,7,8,9}, ymin = 0, ymax=10,
        height = 3in]
       
    \end{axis}
            
    \end{tikzpicture}
    \end{center}}
    
    \begin{solution}
    \begin{center}
         \begin{tikzpicture}
        \begin{axis}[xmin=0, xmax = 10, xtick={0,1,2,3,4,5,6,7,8,9}, ytick={0,1,2,3,4,5,6,7,8,9}, ymin = 0, ymax=10,
        height = 3in]
        \addplot[domain=0:10]{x};
            \addplot[domain=0:10]{3*sqrt(x)};
            \addplot[domain=0:10]{x};
            \addplot[domain=0:10]{3*sqrt(x)};
            \addplot+[domain=0:9, fillregion, draw=none] {3*sqrt(x)} -- (0,
            0);
    \end{axis}
            
    \end{tikzpicture}
    \end{center}
           
    \end{solution}
\item  On your sketch above, slice the region $\mathcal{R}$ vertically into $n$ slices of equal thickness $\Delta x$, and label the $x$-coordinates of the slices as $x_0, x_1, \ldots, x_n$ as usual. 



      \begin{solution}
          See the sketch in the solution to part (e). 
      \end{solution}
      
\item What is $\Delta x$ in terms of $n$? 

\begin{solution}
    We're slicing the interval $[0, 9]$, so $\Delta x = \frac{9 - 0}{n}
      = \frac{9}{n}$.
\end{solution}

\pspace{\vfill}
\item What is $x_k$ in terms
      of $\Delta x$? 
   

    \begin{solution}
      $x_k = 0 + k \Delta x$.
    \end{solution}
\pspace{\vfill}
  \item On your sketch of the region, draw a rectangle that approximates the
      $k$-th slice.  What is the area of this rectangle? Express your answer in terms of $x_k$ and $\Delta x$.
    

    \begin{solution}
    
      \begin{center}
        \begin{tikzpicture}[declare function={f(\x)=sqrt(\x); g(\x) =
            \x/2;}] 
          \begin{axis}[width=3in, height=2in, xtick=\empty, ytick=\empty,
            axis on top=true, ymin=0, axis x line=bottom, xtick={0,
              0.307692, ..., 4.01}, xticklabels={$x_0$, $x_1$,
              $\cdots$,,,,,$x_k$,,,,,,$x_n$}, disabledatascaling]
            \addplot+[domain=0:4, fillregion, draw=none] {sqrt(x)} -- (0,
            0); %
            \addplot+[vertslices] coordinates { (0, {f(0)}) (4/13,
              {f(4/13)}) (8/13, {f(8/13)}) (12/13, {f(12/13)}) (16/13,
              {f(16/13)}) (20/13, {f(20/13)}) (24/13, {f(24/13)}) (28/13,
              {f(28/13)}) (32/13, {f(32/13)}) (36/13, {f(36/13)}) (40/13,
              {f(40/13)}) (44/13, {f(44/13)}) (48/13, {f(48/13)}) }; %
            \addplot+[fill=white, draw=none, domain=0:4] {x/2}
            \closedcycle; %
            \addplot+[domain=0:4.5] {sqrt(x)}; %
            \addplot+[domain=0:4.5] {x/2}; %
            \draw[fill=cyan, fill opacity=0.75, draw=blue] (24/13,
            {f(28/13)}) -- (28/13, {f(28/13)}) -- (28/13, {g(28/13)}) --
            (24/13, {g(28/13)}) -- cycle; %
          \end{axis}
        \end{tikzpicture}
      \end{center}
      This rectangle's width is $\Delta x$, and its height is (top
      $y$-coordinate) $-$ (bottom $y$-coordinate) = $3\sqrt{x_k} -
      x_k$.  Therefore, the rectangle's area is $\boxed{\left(
          3\sqrt{x_k} -
      x_k \right) \Delta x}$.
    \end{solution}
\pspace{\vfill}
  \item
    \begin{problem}
      Write a Riemann sum that approximates the area of $\mathcal{R}$. 
    \end{problem}

    \begin{solution}
      To approximate the area of the whole region $\mathcal{R}$, we just
      add up our approximations for the areas of all the slices, giving
      $\boxed{\sum_{k=1}^n \left( 3\sqrt{x_k} -
      x_k \right)
        \Delta x}$.
    \end{solution}
\pspace{\vfill}

  \item
\begin{problem}
      Take a limit of the Riemann sum in part (f) to write a definite integral that gives the exact area of
      $\mathcal{R}$. Then, evaluate this integral.
    \end{problem}

    \begin{solution}
      Taking the limit as $n \to \infty$ of our Riemann sum, the exact
      area of $\mathcal{R}$ is $\boxed{\int_0^9 \left( 3\sqrt{x} -
      x \right) dx}$. 

      $\displaystyle \int_{0}^{9} (3\sqrt{x}-x)\,dx = \displaystyle \int_{0}^{9} (3x^{1/2}-x)\,dx = \displaystyle \left. \left(2x^{3/2}-\frac{1}{2}x^2\right) \right|_{0}^{9} = 54 - \frac{81}{2} = \frac{27}{2}$
    \end{solution}
 \pspace{\vfill \vfill}
  \end{enumerate}
  

\pspace{\newpage}

\item 

     \begin{enumerate}
     
        \item  Evaluate $\displaystyle \int_0^{\sqrt{3}/2} \arcsin x\,dx$ by visualizing it as an area, and slicing this area horizontally to get an
     equivalent integral which is easier to evaluate. (This area is represented in the sketch below.) 
     
     \begin{framed}
     Note: Horizontal slices have height $\Delta y$ and width (right $x$-coordinate) - (left $x$-coordinate). Remember to rewrite boundaries given in terms of $x$ (i.e., $y=\arcsin x$) into terms of $y$.
     \end{framed}

 \probonly{
\begin{center}
       \begin{tikzpicture}
         \begin{axis}[xtick={0.866}, xlabel={$x$}, xticklabels={$\frac{\sqrt{3}}{2}$},
           ytick={1.047}, ylabel={$y$}, ylabel style={rotate=-90},  yticklabels={}, clip=false,
           width=3in, height=3in]
           \addplot+[domain=0:sqrt(3)/2, fillregion] {rad(asin(x))}
           \closedcycle; %
           \addplot+[domain=-0.2:1.2] ({sin(deg(x))}, x)
          node[above] {$y = \arcsin x$};
         \end{axis}
       \end{tikzpicture}
  \end{center}}

   \begin{solution}
     We can visualize $\displaystyle \int_0^{\sqrt{3}/2} \arcsin x\,dx$ as
     the area under $y = \arcsin x$ from $x = 0$ to $x =
     \frac{\sqrt{3}}{2}$:
     \begin{center}
       \begin{tikzpicture}
         \begin{axis}[xtick={0.866}, xticklabels={$\frac{\sqrt{3}}{2}$},
           ytick={1.047}, yticklabels={$\frac{\pi}{3}$}, clip=false,
           width=2in, height=2.5in]
           \addplot+[domain=0:sqrt(3)/2, fillregion] {rad(asin(x))}
           \closedcycle; %
           \addplot+[domain=-pi/2:pi/2] ({sin(deg(x))}, x)
          node[above] {$y = \arcsin x$};
         \end{axis}
       \end{tikzpicture}
     \end{center}
     (We've labeled $\arcsin \left(\frac{\sqrt{3}}{2} \right) =
     \frac{\pi}{3}$ on the $y$-axis, as that's the highest point in the
     region.)  We're going to slice horizontally, which means we're slicing
     the interval $\left[ 0, \frac{\pi}{3} \right]$ on the $y$-axis into
     $n$ slices of equal height $\Delta y$.  We'll label $y_0 = 0, \ldots,
     y_n = \frac{\pi}{3}$ as usual and focus on the $k$-th slice, which we
     approximate by a rectangle:
     \begin{center}
       \begin{tikzpicture}
         \begin{axis}[xtick={0.866}, xticklabels={$\frac{\sqrt{3}}{2}$},
          ytick={1.047, 0.698}, yticklabels={$\frac{\pi}{3}$, $y_k$}, axis
      on top, disabledatascaling, clip=false, width=2in, height=2.5in]
          \addplot+[domain=0:sqrt(3)/2, fillregion] {rad(asin(x))} \closedcycle;
          \addplot[horizslices] coordinates {
             ({sqrt(3)/2}, 0) ({sqrt(3)/2}, pi/27) ({sqrt(3)/2}, 2*pi/27)
             ({sqrt(3)/2}, 3*pi/27) ({sqrt(3)/2}, 4*pi/27) ({sqrt(3)/2}, 5*pi/27)
             ({sqrt(3)/2}, 6*pi/27) ({sqrt(3)/2}, 7*pi/27) ({sqrt(3)/2}, 8*pi/27)
             ({sqrt(3)/2}, 9*pi/27)};
           \addplot+[domain=0:sqrt(3)/2, fill=white, draw=none]
           {rad(asin(x))} -- (0, {pi/3}) -- (0, 0);
           \addplot+[domain=-pi/2:pi/2] ({sin(deg(x))}, x)
           node[above] {$y = \arcsin x$};
           \draw[fill=cyan, draw=blue, fill opacity=0.75] ({sqrt(3)/2},
          5*pi/27) -- ({sin(deg(6*pi/27))}, 5*pi/27) --
           ({sin(deg(6*pi/27))}, 6*pi/27) -- ({sqrt(3)/2},
           6*pi/27) -- cycle;
           \draw[dashed] (0, 6*pi/27) -- ({sin(deg(6*pi/27))}, 6*pi/27);
         \end{axis}
       \end{tikzpicture}
     \end{center}
     This rectangle's height is $\Delta y$, and its width is (right
     $x$-value) $-$ (left $x$-value).  The right $x$-value is
     $\frac{\sqrt{3}}{2}$; since $y = \arcsin x$ can be rewritten as $x =
     \sin y$, the left $x$-value is $\sin (y_k)$.  Therefore,
     \[ \text{area of $k$-th slice} \approx \left( \frac{\sqrt{3}}{2} -
       \sin y_k \right) \Delta y.\] Summing the areas of all the slices and
     taking the limit as the number of slices $\to \infty$, the area of the
     region is 
     \begin{align*}
       \int_0^{\pi/3} \left( \frac{\sqrt{3}}{2} - \sin y \right) dy &=
       \left. \frac{\sqrt{3}}{2} y + \cos y \right|_0^{\pi/3} \\
       &= \left( \frac{\sqrt{3}}{2} \cdot \frac{\pi}{3} + \frac{1}{2}
       \right) - (0 + 1) \\
       &= \boxed{\frac{\pi \sqrt{3}}{6} - \frac{1}{2}}
     \end{align*}
   \end{solution}
   \pspace{\newpage}
    \item Evaluate $\displaystyle \int_0^{\sqrt{3}/2} \arcsin x\,dx$ using integration by parts. 
    
        \begin{solution}
            Let's evaluate the indefinite integral $\displaystyle \int \arcsin x\,dx$ with integration by parts and then evaluate its difference between $x=0$ and $x=\sqrt{3}/2$.
\begin{center}
    \begin{tabular}{ l l }
  $u=\arcsin x$ &  $dv = 1 dx$  \\
  $du=\dfrac{1}{\sqrt{1-x^2}}dx$ &  $v=x$  \\
\end{tabular}
\end{center}

According to our choices for $u$ and $dv$ above,

\begin{center}
$\displaystyle \int \arcsin x\,dx = x\arcsin x - \displaystyle \int \dfrac{x}{\sqrt{1-x^2}} \,dx$
\end{center}

The remaining integral can be evaluated with the substitution $u=1-x^2$ and $du=-2xdx$, where $-\dfrac{du}{2}=xdx$. The substituted integral takes the form

\begin{center}
    $\displaystyle -\dfrac{1}{2} \int \dfrac{1}{\sqrt{u}}\,du = \displaystyle -\dfrac{1}{2} \int u^{-1/2}\,du = -\dfrac{1}{2}\cdot\dfrac{u^{1/2}}{1/2}+C=-\sqrt{u}+C=-\sqrt{1-x^2}+C$
\end{center}

Let's use our result here to finish up the evaluation of the original integral:

\begin{center}
$\displaystyle \int \arcsin x\,dx = x\arcsin x - \displaystyle \int \dfrac{x}{\sqrt{1-x^2}} \,dx=x\arcsin x + \sqrt{1-x^2}+C$    
\end{center}

Now, let's evaluate the difference in an antiderivative between $x=0$ and $x=\sqrt{3}/2$.

\begin{align*}
\left. \left(x\arcsin x + \sqrt{1-x^2}\right) \right|_{0}^{\sqrt{3}/2} &= \left( \dfrac{\sqrt{3}}{2}\cdot \arcsin \left( \dfrac{\sqrt{3}}{2}\right)+\sqrt{1-\left(\dfrac{\sqrt{3}}{2}\right)^2}\right) - 1\\
&=\dfrac{\sqrt{3}}{2}\cdot\dfrac{\pi}{3}+\dfrac{1}{2}-1\\
&=\dfrac{\pi\sqrt{3}}{6}-\dfrac{1}{2}
\end{align*}
        \end{solution}
        \pspace{\newpage}

    \end{enumerate}
    \end{enumerate}
   \pspace{\newpage}
   

 \pspace{\newpage}
 \htitle{10 Practice Problems for Integration by Parts}
\underline{Directions}: Evaluate the following indefinite integrals.
\begin{enumerate}
\item $\displaystyle \int x^2\ln x \,dx$

\begin{solution}
\begin{center}
    \begin{tabular}{ l l }
  $u=\ln x$ &  $dv = x^2 dx$  \\
  $du=\frac{1}{x}dx$ &  $v=\frac{1}{3}x^3$  \\
\end{tabular}
\begin{align*}
\displaystyle \int x^2\ln x \,dx&=\frac{1}{3}x^3\ln x - \displaystyle \dfrac{1}{3}\int x^2 \,dx \\&= \frac{1}{3}x^3\ln x - \dfrac{1}{9}x^3+C
\end{align*}
\end{center} 
\end{solution}
\vfill
\item $\displaystyle \int x^2e^{2x+1}\,dx$

\begin{solution}
\begin{center}
    \begin{tabular}{ l l }
  $u=x^2$ &  $dv = e^{2x+1} dx$  \\
  $du=2xdx$ &  $v=\frac{1}{2}e^{2x+1}$  \\
\end{tabular}
 
$\displaystyle \int x^2e^{2x+1}\,dx=\dfrac{1}{2}x^2e^{2x+1}-\displaystyle \int xe^{2x+1}\,dx$
\end{center} 
\begin{center}
    \begin{tabular}{ l l }
  $u=x$ &  $dv = e^{2x+1} dx$  \\
  $du=1dx$ &  $v=\frac{1}{2}e^{2x+1}$  \\
\end{tabular}

\begin{align*}
\displaystyle \int x^2e^{2x+1}\,dx&=\dfrac{1}{2}x^2e^{2x+1}-\left[\dfrac{1}{2}xe^{2x+1}-\dfrac{1}{2}\int e^{2x+1}\,dx\right]\\
&=\dfrac{1}{2}x^2e^{2x+1}+\dfrac{1}{2}xe^{2x+1}+\dfrac{1}{4}e^{2x+1}+C
\end{align*}
\end{center}
\end{solution}
\vfill

\item $\displaystyle \int x\ln(1+x)\,dx$

\begin{solution}
First, let's use the substitution $y=1+x$.
\[
    \int x\ln(1+x)\,dx 
    = \int (y-1)\ln y\, dy
\]
    Now we can use integration by parts with
\begin{center}
\begin{tabular}{ l l }
  $u=\ln y$ &  $dv = y-1$ \\
  $du=\frac{1}{y}dy$ &  $v=\frac{1}{2}y^2-y$  \\
\end{tabular}
\begin{align*}
   \int (y-1)\ln y\, dy
    &= (\tfrac{1}{2}y^2-y)\ln y
        - \int \frac{\frac{1}{2}y^2-y}{y}\, dy\\
    &= (\tfrac{1}{2}y^2-y)\ln y - \int (\tfrac{1}{2}y-1)\,dy\\
    &= (\tfrac{1}{2}y^2-y)\ln y
        - (\tfrac{1}{4}y^2 - y) +C\\
    &= \tfrac{1}{2}y^2 \ln y - y \ln y - \tfrac{1}{4} y^2 + y + C \\
    &= \frac{1}{2}(1+x)^2\ln(1+x) - (1+x)\ln (x+1) - \frac{1}{4}(1+x)^2 + (1+x) +C
\end{align*}
\end{center}
\end{solution}
\pspace{\vfill}

\item $\displaystyle \int x\sec^2 x\,dx$

\begin{solution}
\begin{center}
    \begin{tabular}{ l l }
  $u=x$ &  $dv = \sec^2 x dx$  \\
  $du=1dx$ &  $v=
  \tan x$  \\
\end{tabular}
\begin{align*}
    \int x\sec^2 x\,dx &= x\tan x - \int \tan x\,dx \\
    &= x\tan x + \ln|\cos x| + C \\
\end{align*}
\end{center}\vspace{-24pt}
Remember that we previously derived $\int \tan x \,dx = -\ln|\cos x| +C$ using substitution.
\end{solution}
\pspace{\vfill}
\pspace{\newpage}
\item $\displaystyle \int \cos x \ln(\sin x)\,dx$

\begin{solution}
First, let's use the substitution $u = \sin x$, $du = \cos x\, dx$.
\[\int \cos x \ln(\sin x)\, dx = \int \ln u \, du\]
Now we have $\int \ln u \, du$, which is $u\ln u - u +C$, as we've derived previously using integration by parts. 
\begin{align*}
    \int \cos x \ln(\sin x)\,dx &= \int \ln u \,du \\
    &= u\ln u - u + C \\
    &= (\sin x) \ln(\sin x) - \sin x + C
\end{align*}
\end{solution}
\pspace{\vfill}

\item $\displaystyle \int \ln(x^2+1)\,dx$ 
\begin{framed}
    During an intermediate step, the following fact may be helpful: 
    \begin{center}
        $\dfrac{x^2}{x^2+1}=\dfrac{(x^2+1)-1}{x^2+1}$
    \end{center}
\end{framed}

\begin{solution}
Let's use integration by parts.
\begin{center}
    \begin{tabular}{ l l }
  $u=\ln(x^2+1)$ &  $dv =1 dx$  \\
  $du=\frac{2x}{x^2+1}dx$ &  $v=x$  \\
\end{tabular}
\end{center} \vspace{-12pt}
\begin{align*}
    \int \ln(x^2+1)\,dx &= x\ln(x^2+1) - \int \frac{2x^2}{x^2+1}\,dx \\
    \intertext{To do the remaining integral, we can rewrite it as follows.}
    &= x\ln(x^2+1) - 2\int \frac{x^2}{x^2+1}\,dx \\
    &= x\ln(x^2+1) - 2\int \frac{x^2+1-1}{x^2+1}\,dx \\
    &= x\ln(x^2+1) - 2\left(\int \frac{x^2+1}{x^2+1}\,dx
        - \int \frac{1}{x^2+1}\,dx\right)\\
    &= x\ln(x^2+1) - 2\left(\int dx - \int \frac{1}{x^2+1}\,dx \right)\\
    &= x\ln(x^2+1) - 2(x - \arctan x) + C
\end{align*}
   
\end{solution}



\pspace{\vfill}

\item $\displaystyle \int \arctan x\,dx$ 

\begin{solution}
\begin{center}
    \begin{tabular}{ l l }
  $u=\arctan x$ &  $dv =1 dx$  \\
  $du=\frac{1}{1+x^2}dx$ &  $v=x$  \\
\end{tabular}

$\displaystyle \int \arctan x\,dx = x\arctan x - \displaystyle \int \dfrac{x}{1+x^2}\,dx$
\end{center}
Let's start to solve the remaining integral with the substitution $u=1+x^2$ and $du=2xdx$, where $\dfrac{du}{2}=xdx$.
\begin{center}
$\displaystyle \int \dfrac{x}{1+x^2}\,dx=\dfrac{1}{2}\displaystyle \int \dfrac{1}{u}\,du=\dfrac{1}{2}\ln |u|+C=\dfrac{1}{2}\ln(1+x^2)+C$ 
\end{center}
Thus,
\begin{center}
$\displaystyle \int \arctan x\,dx = x\arctan x - \dfrac{1}{2}\ln(1+x^2)+C$    
\end{center}
\end{solution}
\pspace{\vfill}
\pspace{\newpage}
\item $\displaystyle \int \sin x \sec x \tan x\,dx$ 

\begin{solution}
    \begin{center}
    \begin{tabular}{ l l }
  $u=\sin x$ &  $dv =\sec x \tan x dx$  \\
  $du=\cos xdx$ &  $v=\sec x$  \\
\end{tabular}
\begin{align*}
\displaystyle \int \sin x \sec x \tan x\,dx &= \sin x \sec x - \displaystyle \int \sec x \cos x \,dx \\ &= \sin x \sec x - \displaystyle \int 1 \,dx \\ &= \sin x \sec x - x + C \\ &= \tan x - x + C
\end{align*}
\end{center}
\end{solution}
\pspace{\vfill}
\item $\displaystyle \int (\ln x)^2\,dx$

\begin{solution}
\begin{center}
    \begin{tabular}{ l l }
  $u=(\ln x)^2$ &  $dv = 1 dx$  \\
  $du=\frac{2\ln x}{x}dx$ &  $v=x$  \\
\end{tabular}
\begin{align*}
\displaystyle \int (\ln x)^2\,dx &= x(\ln x)^2 - \displaystyle 2 \int \ln x\,dx     
\end{align*}
\begin{tabular}{ l l }
  $u=\ln x$ &  $dv = 1 dx$  \\
  $du=\frac{1}{x}dx$ &  $v=x$  \\
\end{tabular}
\begin{align*}
   \displaystyle \int (\ln x)^2\,dx &= x(\ln x)^2 - \displaystyle 2\left [x \ln x - \displaystyle \int 1\,dx \right]\\
   &= x(\ln x)^2 - 2x \ln x +2x +C
\end{align*}
\end{center}    
\end{solution}
\pspace{\vfill}

\item $\displaystyle \int (\ln x)^3\,dx$ (Note: This problem is exploratory, and very challenging! Before starting, be sure to recall $\int \ln x \,dx$, and also confirm your answer to $\int (\ln x)^2 \,dx$ so you can let $dv=(\ln x)^2dx$.)

\begin{solution}
\begin{center}
\begin{tabular}{ l l }
  $u=\ln x$ &  $dv = (\ln x)^2 dx$  \\
  $du=\frac{1}{x}dx$ &  $v=x(\ln x)^2 - 2x \ln x +2x$  \\
\end{tabular}
\begin{align*}
\displaystyle \int (\ln x)^3\,dx&=\ln x \left[ x(\ln x)^2 - 2x \ln x +2x \right] - \displaystyle \int \dfrac{x(\ln x)^2 - 2x \ln x +2x}{x}\,dx \\ &=x(\ln x)^3-2x (\ln x)^2+2x\ln x - \displaystyle \int \left[ (\ln x)^2 - 2 \ln x +2 \right]\,dx \\ 
&= x(\ln x)^3-2x (\ln x)^2+2x\ln x - \left[x(\ln x)^2 - 2x \ln x +2x -2x\ln x +2x +2x\right]+C\\
&= x(\ln x)^3-2x (\ln x)^2+2x\ln x - x(\ln x)^2 +2x\ln x-2x + 2x\ln x-2x-2x 
\\ &= x(\ln x)^3-3x (\ln x)^2+6x\ln x-6x+C
\end{align*}
\end{center}

\end{solution}
\pspace{\vfill}
\end{enumerate}

\end{document}